\documentclass[12pt]{article}

\usepackage{sbc-template}
\usepackage{graphicx,url}
\usepackage[utf8]{inputenc}
\usepackage[T1]{fontenc}
\usepackage[brazil]{babel}
\usepackage{booktabs}
\usepackage{amsmath}
\usepackage{array}

\graphicspath{{../../results/figures/}{results/figures/}}

\sloppy

\title{Sistema de Recomenda\c{c}\~ao de Produtos Utilizando Apache Spark}

\author{Nome do Autor\inst{1}}

\address{
  Instituto de Computa\c{c}\~ao -- Universidade Federal de Alagoas (UFAL)\\
  Macei\'o -- AL -- Brasil
  \email{email@exemplo.com}
}

\begin{document}

\maketitle

\begin{abstract}
This paper presents a local big data pipeline for product recommendation
experiments using Apache Spark and the RetailRocket e-commerce dataset. The
pipeline covers data ingestion, cleaning, implicit feedback weighting,
temporal splitting, recommendation model training, top-$k$ evaluation and
runtime benchmarking. Three recommendation strategies are compared:
global popularity, item-item co-occurrence and implicit-feedback ALS. The
full benchmark shows that the co-occurrence model achieved the best
Precision@10 and Recall@10 in the evaluated setting, while ALS presented
lower catalog coverage and higher sensitivity to hyperparameters. The
results provide empirical evidence about quality, scalability,
partitioning and storage trade-offs in a reproducible local Spark
environment.
\end{abstract}

\begin{resumo}
Este artigo apresenta uma pipeline local de dados massivos para experimentos
de recomenda\c{c}\~ao de produtos utilizando Apache Spark e o conjunto de
dados RetailRocket. A pipeline contempla ingest\~ao, limpeza, pondera\c{c}\~ao
de feedback impl\'icito, divis\~ao temporal, treinamento de modelos de
recomenda\c{c}\~ao, avalia\c{c}\~ao top-$k$ e medi\c{c}\~ao de tempo de
execu\c{c}\~ao. S\~ao comparadas tr\^es estrat\'egias: popularidade global,
coocorr\^encia item-item e ALS para feedback impl\'icito. O benchmark completo
indica que o modelo de coocorr\^encia obteve os melhores valores de
Precision@10 e Recall@10 no ambiente avaliado, enquanto o ALS apresentou
menor cobertura de cat\'alogo e maior sensibilidade aos hiperpar\^ametros. Os
resultados oferecem evid\^encias sobre qualidade, escalabilidade,
particionamento e armazenamento em uma execu\c{c}\~ao Spark local e
reprodut\'ivel.
\end{resumo}

\section{Introdu\c{c}\~ao}

Sistemas de recomenda\c{c}\~ao s\~ao componentes centrais em aplica\c{c}\~oes de
com\'ercio eletr\^onico, pois auxiliam usu\'arios a encontrar produtos em
cat\'alogos extensos e din\^amicos. Nesses cen\'arios, os sinais dispon\'iveis
frequentemente s\~ao impl\'icitos, como visualiza\c{c}\~oes, adi\c{c}\~oes ao
carrinho e compras, em vez de avalia\c{c}\~oes expl\'icitas. Al\'em disso, a
quantidade de eventos e a alta esparsidade das intera\c{c}\~oes tornam
necess\'ario o uso de ferramentas adequadas para processamento distribu\'ido ou
paralelo.

Apache Spark \cite{zaharia2016spark} oferece uma plataforma unificada para
processamento de dados em larga escala, com suporte a transforma\c{c}\~oes,
agrega\c{c}\~oes e algoritmos de aprendizado de m\'aquina por meio da MLlib.
Este trabalho investiga a constru\c{c}\~ao de uma pipeline local e
reprodut\'ivel para benchmark de recomenda\c{c}\~ao de produtos com Spark,
utilizando o conjunto RetailRocket \cite{retailrocket2015dataset}.

O objetivo \'e avaliar, em um ambiente local, o comportamento de modelos
simples e escal\'aveis de recomenda\c{c}\~ao em termos de qualidade de ranking,
tempo de execu\c{c}\~ao, impacto do particionamento e formato de armazenamento.
O trabalho privilegia simplicidade e reprodutibilidade, evitando servi\c{c}os
externos e mantendo todas as etapas principais acess\'iveis por linha de
comando.

\section{Fundamenta\c{c}\~ao}

\subsection{Feedback impl\'icito}

Em feedback impl\'icito, a prefer\^encia do usu\'ario \'e inferida a partir de
eventos observados. Neste trabalho, eventos de visualiza\c{c}\~ao,
adi\c{c}\~ao ao carrinho e compra recebem pesos distintos, respectivamente
$1$, $3$ e $5$. A intera\c{c}\~ao final entre usu\'ario e produto \'e obtida por
agrega\c{c}\~ao dos pesos e, por padr\~ao, transformada por $\log(1+x)$ para
reduzir a influ\^encia de valores extremos.

\subsection{Modelos avaliados}

O modelo de popularidade recomenda os produtos com maior escore agregado no
conjunto de treino. Apesar de simples, ele \'e uma linha de base importante em
recomenda\c{c}\~ao, especialmente em cen\'arios de alta esparsidade.

O modelo de coocorr\^encia item-item utiliza os hist\'oricos de usu\'arios para
estimar associa\c{c}\~oes entre produtos. Produtos vistos em conjunto por muitos
usu\'arios recebem maior similaridade e s\~ao utilizados para gerar
recomenda\c{c}\~oes a partir dos itens j\'a observados no treino.

O ALS (\textit{Alternating Least Squares}) para feedback impl\'icito
\cite{hu2008collaborative} fatora a matriz usu\'ario-item em vetores latentes.
Neste trabalho foram testadas combina\c{c}\~oes de posto, regulariza\c{c}\~ao e
par\^ametro de confian\c{c}a definidas na configura\c{c}\~ao da pipeline.

\subsection{M\'etricas}

A avalia\c{c}\~ao utiliza m\'etricas top-$k$ com $k=10$: Precision@10,
Recall@10, MAP@10, NDCG@10 e cobertura. Precision@10 mede a fra\c{c}\~ao da
lista recomendada que \'e relevante; Recall@10 mede a fra\c{c}\~ao dos itens
relevantes recuperada; MAP@10 e NDCG@10 consideram a posi\c{c}\~ao dos acertos
no ranking; cobertura mede a propor\c{c}\~ao do cat\'alogo que aparece nas
recomenda\c{c}\~oes.

\section{Metodologia}

\subsection{Dataset}

Foi utilizado o RetailRocket eCommerce Dataset, disponibilizado no Kaggle
\cite{retailrocket2015dataset}. O arquivo principal de eventos cont\'em
intera\c{c}\~oes de usu\'arios com produtos, incluindo visualiza\c{c}\~oes,
adi\c{c}\~oes ao carrinho e transa\c{c}\~oes. A Tabela~\ref{tab:dataset}
apresenta o perfil ap\'os a etapa de limpeza.

\begin{table}[!htbp]
\centering
\caption{Perfil do conjunto de dados ap\'os limpeza.}
\label{tab:dataset}
\begin{tabular}{lr}
\toprule
Caracter\'istica & Valor \\
\midrule
Eventos & 2.755.641 \\
Usu\'arios distintos & 1.407.580 \\
Itens distintos & 235.061 \\
Visualiza\c{c}\~oes & 2.664.218 \\
Adi\c{c}\~oes ao carrinho & 68.966 \\
Transa\c{c}\~oes & 22.457 \\
Per\'iodo & 2015-05-03 a 2015-09-17 \\
Esparsidade estimada & 0{,}999992 \\
Eventos por usu\'ario, em m\'edia & 1{,}96 \\
Eventos por item, em m\'edia & 11{,}72 \\
\bottomrule
\end{tabular}
\end{table}

\subsection{Pipeline experimental}

A pipeline foi implementada em Python com PySpark. Inicialmente, os arquivos
CSV s\~ao lidos com esquemas expl\'icitos. Em seguida, eventos com campos
essenciais ausentes s\~ao removidos, eventos inv\'alidos s\~ao filtrados e
duplicatas s\~ao descartadas. Os timestamps s\~ao convertidos para data e hora,
e cada evento recebe um peso de feedback impl\'icito.

As intera\c{c}\~oes s\~ao agregadas por par usu\'ario-produto, gerando escore,
contagem de eventos e data da \'ultima intera\c{c}\~ao. O particionamento de
treino e teste segue uma estrat\'egia temporal: aproximadamente 80\% dos
eventos mais antigos s\~ao usados para treino, e os eventos restantes para
teste. A relev\^ancia no teste prioriza eventos de adi\c{c}\~ao ao carrinho e
compra; em amostras sem tais eventos, a pipeline usa visualiza\c{c}\~oes como
fallback.

\subsection{Ambiente de execu\c{c}\~ao}

Os experimentos foram executados em modo local com Spark \texttt{local[*]}.
O ambiente utilizado possui processador Intel Core i5-12450HX, 12 threads
dispon\'iveis, 7{,}6 GiB de mem\'oria RAM, OpenJDK 17 e PySpark 4.1.2. A
configura\c{c}\~ao da aplica\c{c}\~ao definiu 4 GiB de mem\'oria para o driver e
parti\c{c}\~oes de shuffle adaptativas.

\section{Resultados}

\subsection{Compara\c{c}\~ao de modelos}

A Tabela~\ref{tab:modelos} apresenta um resumo por fam\'ilia de modelo no
benchmark completo, combinando as melhores m\'etricas observadas com o tempo
m\'edio das execu\c{c}\~oes daquela fam\'ilia. O modelo de coocorr\^encia obteve a
melhor qualidade em todas as m\'etricas principais e tamb\'em a maior cobertura
de cat\'alogo. Os modelos ALS apresentaram desempenho superior ao baseline de
popularidade em algumas configura\c{c}\~oes, mas com cobertura baixa em rela\c{c}\~ao
ao cat\'alogo total.

\begin{table}[!htbp]
\centering
\caption{Melhores resultados por fam\'ilia de modelo no benchmark completo.}
\label{tab:modelos}
\begin{tabular}{lrrrrrr}
\toprule
Modelo & Prec@10 & Rec@10 & MAP@10 & NDCG@10 & Cob. & Tempo m\'edio (s) \\
\midrule
Coocorr\^encia & 0{,}012891 & 0{,}054153 & 0{,}022496 & 0{,}034383 & 0{,}700708 & 35{,}56 \\
ALS rank 16 & 0{,}004687 & 0{,}019536 & 0{,}005948 & 0{,}011882 & 0{,}004671 & 11{,}56 \\
ALS rank 8 & 0{,}004175 & 0{,}011999 & 0{,}005199 & 0{,}008014 & 0{,}002712 & 7{,}94 \\
Popularidade & 0{,}002584 & 0{,}010390 & 0{,}003267 & 0{,}006179 & 0{,}000409 & 10{,}18 \\
\bottomrule
\end{tabular}
\end{table}

A Figura~\ref{fig:precision-recall} resume a compara\c{c}\~ao entre
Precision@10 e Recall@10 por modelo e configura\c{c}\~ao, enquanto a
Figura~\ref{fig:coverage} apresenta a cobertura.

\begin{figure}[!htbp]
\centering
\includegraphics[width=.9\textwidth]{precision_recall_by_model.png}
\caption{Precision@10 e Recall@10 por modelo avaliado.}
\label{fig:precision-recall}
\end{figure}

\begin{figure}[!htbp]
\centering
\includegraphics[width=.78\textwidth]{coverage_by_model.png}
\caption{Cobertura de cat\'alogo por modelo avaliado.}
\label{fig:coverage}
\end{figure}
\clearpage

\subsection{Escalabilidade por fra\c{c}\~ao do dataset}

O benchmark de escala foi executado com fra\c{c}\~oes de 25\%, 50\%, 75\% e
100\% dos usu\'arios, usando amostragem determin\'istica por identificador de
usu\'ario. A Tabela~\ref{tab:escala} apresenta os tempos totais registrados
para cada fra\c{c}\~ao. Esses tempos representam a execu\c{c}\~ao do conjunto de
modelos avaliado em cada fra\c{c}\~ao. Assim, a interpreta\c{c}\~ao desta tabela se
concentra no volume processado e no custo temporal; a compara\c{c}\~ao de
qualidade entre modelos \'e apresentada na Tabela~\ref{tab:modelos}. A
Figura~\ref{fig:runtime-fraction} mostra a varia\c{c}\~ao do tempo total ao
aumentar a fra\c{c}\~ao processada.

\begin{table}[!htbp]
\centering
\caption{Benchmark por fra\c{c}\~ao do dataset.}
\label{tab:escala}
\begin{tabular}{rrrrr}
\toprule
Fra\c{c}\~ao & Eventos de entrada & Usu\'arios & Itens & Tempo total (s) \\
\midrule
0{,}25 & 677.301 & 352.154 & 134.181 & 106{,}03 \\
0{,}50 & 1.370.805 & 704.673 & 181.066 & 139{,}92 \\
0{,}75 & 2.061.925 & 1.055.538 & 212.192 & 118{,}20 \\
1{,}00 & 2.755.641 & 1.407.580 & 235.061 & 193{,}55 \\
\bottomrule
\end{tabular}
\end{table}

\begin{figure}[!htbp]
\centering
\includegraphics[width=.78\textwidth]{runtime_by_fraction.png}
\caption{Tempo total por fra\c{c}\~ao do dataset.}
\label{fig:runtime-fraction}
\end{figure}
\clearpage

\subsection{Particionamento}

A Tabela~\ref{tab:particionamento} mostra o tempo total observado para 2, 4, 8
e 16 parti\c{c}\~oes na fra\c{c}\~ao de 25\% do dataset, usando o modelo de
popularidade como refer\^encia. O melhor tempo local foi obtido com 8
parti\c{c}\~oes, embora 4 parti\c{c}\~oes tenha apresentado valor pr\'oximo. A
Figura~\ref{fig:runtime-partitions} evidencia que o aumento de parti\c{c}\~oes
nem sempre reduz o tempo total em execu\c{c}\~ao local.

\begin{table}[!htbp]
\centering
\caption{Benchmark de particionamento com 25\% do dataset.}
\label{tab:particionamento}
\begin{tabular}{rrrrr}
\toprule
Parti\c{c}\~oes & Tempo total (s) & Prec@10 & Rec@10 & MAP@10 \\
\midrule
2 & 14{,}66 & 0{,}001709 & 0{,}009402 & 0{,}002422 \\
4 & 12{,}34 & 0{,}001724 & 0{,}009483 & 0{,}002443 \\
8 & 12{,}13 & 0{,}001724 & 0{,}009483 & 0{,}002443 \\
16 & 14{,}97 & 0{,}001709 & 0{,}009402 & 0{,}002422 \\
\bottomrule
\end{tabular}
\end{table}

\begin{figure}[!htbp]
\centering
\includegraphics[width=.78\textwidth]{runtime_by_partitions.png}
\caption{Tempo total por n\'umero de parti\c{c}\~oes.}
\label{fig:runtime-partitions}
\end{figure}
\clearpage

\subsection{CSV e Parquet}

O benchmark de armazenamento comparou leitura e contagem do CSV original com
uma vers\~ao Parquet derivada dos eventos limpos. O Parquet reduziu o tamanho
em disco de 89{,}87 MiB para 57{,}34 MiB. Nesta execu\c{c}\~ao local, por\'em, o
tempo de leitura e contagem foi ligeiramente menor para o CSV: 0{,}53 s contra
0{,}60 s no Parquet. Esse resultado sugere que, para uma opera\c{c}\~ao simples
de contagem em ambiente local, o ganho de armazenamento n\~ao necessariamente
se traduz em menor lat\^encia; em consultas seletivas e pipelines repetidas, o
formato colunar tende a ser mais vantajoso.
A Figura~\ref{fig:storage} sintetiza a compara\c{c}\~ao de tempo entre os dois
formatos.

\begin{figure}[!htbp]
\centering
\includegraphics[width=.75\textwidth]{storage_csv_vs_parquet.png}
\caption{Tempo de leitura e contagem para CSV e Parquet.}
\label{fig:storage}
\end{figure}
\clearpage

\section{Discuss\~ao}

Os resultados indicam que, para o conjunto de dados e a metodologia adotados,
a coocorr\^encia item-item foi a alternativa mais competitiva. Uma hip\'otese
para esse comportamento \'e que o dataset apresenta intera\c{c}\~oes muito
esparsas, com m\'edia inferior a dois eventos por usu\'ario. Nesse cen\'ario,
modelos latentes como ALS podem ter dificuldade para aprender representa\c{c}\~oes
est\'aveis para muitos usu\'arios e itens, enquanto regras de coocorr\^encia
capturam associa\c{c}\~oes locais de maneira direta.

A baixa cobertura dos modelos ALS tamb\'em merece aten\c{c}\~ao. Como as
recomenda\c{c}\~oes foram geradas apenas para usu\'arios avali\'aveis e itens
vistos no treino foram removidos, a lista final pode ficar concentrada em uma
fra\c{c}\~ao pequena do cat\'alogo. J\'a a coocorr\^encia apresentou cobertura
superior a 65\% nas melhores execu\c{c}\~oes, o que sugere maior diversidade de
itens recomendados.

Os tempos observados devem ser interpretados como evid\^encias de execu\c{c}\~ao
local, n\~ao como medi\c{c}\~oes de cluster. A execu\c{c}\~ao em uma \'unica m\'aquina
sofre influ\^encia de cache, escalonamento local, parti\c{c}\~oes e custo de
a\c{c}\~oes Spark acionadas durante a avalia\c{c}\~ao.

\section{Amea\c{c}as \`a Validade}

Este estudo utiliza avalia\c{c}\~ao offline e, portanto, n\~ao mede impacto real
em convers\~ao, receita ou satisfa\c{c}\~ao do usu\'ario. O dataset \'e anonimizado
e n\~ao possui sess\~oes expl\'icitas, o que limita a modelagem sequencial. Al\'em
disso, o ambiente local n\~ao representa um cluster Spark de produ\c{c}\~ao. Por
fim, a grade de hiperpar\^ametros do ALS foi pequena, privilegiando custo
computacional controlado em vez de busca extensiva.

\section{Conclus\~ao}

Este artigo apresentou uma pipeline local e reprodut\'ivel para benchmark de
recomenda\c{c}\~ao de produtos com Apache Spark. A partir do dataset RetailRocket,
foram implementadas etapas de limpeza, pondera\c{c}\~ao de eventos,
agrega\c{c}\~ao, split temporal, treino, avalia\c{c}\~ao top-$k$ e an\'alise de
desempenho. Nos experimentos realizados, a coocorr\^encia item-item apresentou
o melhor desempenho de qualidade e a maior cobertura, superando popularidade
e ALS nas configura\c{c}\~oes avaliadas.

Como trabalhos futuros, prop\~oe-se ampliar a busca de hiperpar\^ametros,
avaliar modelos sequenciais, incorporar atributos de produto, executar os
experimentos em cluster Spark e realizar an\'alises estat\'isticas com m\'ultiplas
rodadas para reduzir a influ\^encia de variabilidade local.

\bibliographystyle{sbc}
\bibliography{referencias}

\end{document}
